\chapter{L'idea in parole semplici}\label{cap:p1-idea}

% Capitolo della Parte I --- spiegazione per chi non conosce l'IA.

Immagina di dover guardare sei numeri che descrivono lo stato di una persona
--- la pressione, i battiti del cuore, la temperatura, l'ossigeno nel sangue,
lo zucchero nel sangue --- e di dover rispondere a una sola domanda:
\emph{``questa persona è a rischio basso, medio o alto?''}.

È tutto qui. Il lavoro del sistema si riduce a questo: legge sei numeri e
restituisce una parola tra \classe{basso}, \classe{medio} e \classe{alto}.

\section{Perché servono degli ``aiutanti''}\label{sec:p1-perche-aiutanti}

A prima vista sembra facile: basta scrivere ``se la pressione è troppo alta,
allora rischio alto''. E in effetti una regola del genere si può scrivere. Ma
il corpo umano è complicato: i sei numeri si influenzano a vicenda, e decidere
il rischio giusto richiede di considerarli tutti insieme. Scrivere a mano
tutte le combinazioni possibili è noioso, rigido e facile sbagliare.

Per questo il progetto usa l'intelligenza artificiale: invece di elencare
ogni singola regola, si fa in modo che una macchina \emph{impari} il pattern
dai dati. Ma perché la macchina impari, e perché lo faccia in sicurezza, il
sistema si appoggia a diversi ``aiutanti'', ognuno con un compito preciso.

\begin{itemize}
  \item \textbf{L'esperto (teacher rule-based)}: conosce le regole cliniche e
        sa dire, per ogni combinazione di valori, qual è il rischio giusto.
  \item \textbf{Gli esempi (dataset)}: tanti casi di pazienti ``di prova'',
        ognuno già etichettato con il rischio corretto dall'esperto.
  \item \textbf{Lo studente (MLP, la rete neurale)}: studia quegli esempi e
        impara a riconoscere il pattern, così da giudicare anche pazienti mai
        visti.
  \item \textbf{L'ispettore di sicurezza (safety gate)}: controlla per primo i
        casi gravi e ha l'ultima parola.
  \item \textbf{Il pulsante ``non sono sicuro'' (fallback)}: quando lo studente
        ha dei dubbi, si chiede consiglio all'esperto.
\end{itemize}

I prossimi capitoli spiegano ciascun aiutante con parole semplici. Nella
Parte II troverai la stessa storia, ma dal punto di vista tecnico: dati,
algoritmi, codice e numeri.
